\subsection{Computational validation and transfer protocol}

The validation target is the declared simulator operating envelope, not an
implicit claim of completed field flight. Four checks are applied before a result
enters the manuscript. First, numerical reproducibility is enforced by the
SHA-256 evidence manifest, deterministic seeds and a clean-room rerun of the
legacy receipt. Second, internal system validity is established with paired
aware--agnostic episodes: the simulator, scene, horizon, controller, allocator
interface and endpoint are held fixed while one information or geometry factor
is opened at a time. Third, endpoint validity is checked by reading thresholded
completion together with closest approach and by repeating the force ladder over
capture radii. Fourth, transfer validity is represented by a declared parameter
correspondence rather than by an unreported hardware claim.

The correspondence is operational. A target vehicle is specified by mass,
available thrust, controller update rate, wind-estimator update rate and capture
criterion; the simulator quantities in the parameter contract are the initial
values for those fields. A hardware or hardware-in-the-loop extension should
re-estimate mass and thrust, time-align measured wind with vehicle state, replay
the same constant and temporal profiles at the declared update rate, and run
both arms on identical scenes or recorded traces. The primary acceptance
quantity remains the paired completion contrast at the transition, with the
continuous closest-approach endpoint and the clearance diagnostic reported
alongside it. This protocol makes the current computation a reproducible
pre-flight screening and mechanism-isolation study while leaving physical
transfer as a separately testable stage.

The distinction is part of the contribution. The present paper establishes what
can be measured, reproduced and stress-tested before hardware is committed; it
does not substitute simulator evidence for field evidence. All conclusions are
therefore indexed to the stated rigid-body simulator, force profiles, geometry
cells and homogeneous team sizes, and the transfer protocol identifies exactly
which quantities must be revalidated when those conditions change.
